\documentclass[11pt,a4paper]{article}
\usepackage[utf8]{inputenc}
\usepackage[T1]{fontenc}
\usepackage[catalan]{babel}
\usepackage{lmodern}
\usepackage[sfdefault,lining]{FiraSans}
\usepackage[a4paper,top=2.3cm,bottom=2.2cm,left=1.8cm,right=1.8cm,headheight=1.3cm,headsep=0.4cm,footskip=1cm]{geometry}
\usepackage{amsmath,amssymb}
\usepackage{icomma}
\usepackage{graphicx}
\usepackage[table]{xcolor}
\usepackage{tikz}
\usepackage{fontawesome5}
\usepackage{ccicons}
\usepackage{fancyhdr}
\usepackage{lastpage}
\usepackage{tasks}
\usepackage{needspace}
\usepackage[most]{tcolorbox}
\usepackage{array}
\usepackage[hidelinks,pdftitle={Fitxa d'exercicis de la unitat 4: Àlgebra. Polinomis. 3r d'ESO},pdfsubject={Matemàtiques 3r ESO},pdfkeywords={llenguatge algebraic, monomis, polinomis, suma, resta, producte, identitats notables, factor comú, divisió de polinomis, regla de Ruffini}]{hyperref}

\definecolor{navy}{HTML}{1B2748}
\definecolor{accent}{HTML}{E8603A}
\definecolor{teal}{HTML}{1F9E90}
\definecolor{soft}{HTML}{EEF1F8}
\definecolor{muted}{HTML}{5D6685}

\setlength{\parindent}{0pt}
\setlength{\parskip}{3pt}
\settasks{label=\alph*), label-format=\bfseries\color{navy}, label-width=1.3em, item-indent=1.8em, column-sep=0.6em, after-item-skip=5pt, before-skip=3pt, after-skip=4pt}

% ---------- capçalera i peu ----------
\newcommand{\llicencia}{\href{https://creativecommons.org/licenses/by-nc-sa/4.0/deed.ca}{\ccbyncsa}\ \ Llicència Creative Commons Reconeixement-NoComercial-CompartirIgual 4.0 Internacional (CC BY-NC-SA 4.0)}
\pagestyle{fancy}
\fancyhf{}
\fancyhead[L]{\raisebox{-0.35\height}{\includegraphics[height=0.95cm]{logo-gva.png}}}
\fancyhead[C]{\small\color{muted}Fitxa d'exercicis · Unitat 4: Àlgebra · 3r d'ESO}
\fancyhead[R]{\raisebox{-0.35\height}{\includegraphics[height=0.95cm]{logo-vermellar.png}}}
\fancyfoot[L]{\scriptsize\color{muted}\llicencia}
\fancyfoot[R]{\small\color{muted}\thepage}
\renewcommand{\headrulewidth}{0.4pt}
\renewcommand{\headrule}{\color{muted!40}\hrule height \headrulewidth}
\fancypagestyle{primera}{\fancyhf{}\renewcommand{\headrulewidth}{0pt}
  \fancyfoot[L]{\scriptsize\color{muted}\llicencia}
  \fancyfoot[R]{\small\color{muted}\thepage}}

% ---------- nivells i exercicis ----------
\newcommand{\estrelles}[1]{{\color{accent}\ifcase#1\or\faStar\faStar[regular]\faStar[regular]\or\faStar\faStar\faStar[regular]\or\faStar\faStar\faStar\fi}}
\newcounter{exer}
\newcommand{\ex}[1]{\par\needspace{4\baselineskip}\medskip\stepcounter{exer}\noindent{\color{navy}\bfseries\theexer.}\ \ {\small\estrelles{#1}}\ \ }

\newcommand{\seccio}[2]{\par\needspace{10\baselineskip}\bigskip
  \begin{tcolorbox}[colback=navy,colframe=navy,coltext=white,arc=2mm,boxrule=0pt,left=3mm,right=3mm,top=1.2mm,bottom=1.2mm]
  {\large\bfseries #1\quad #2}
  \end{tcolorbox}\vspace{-2pt}}

\newcolumntype{C}[1]{>{\centering\arraybackslash}p{#1}}

% quadre de recordatori (només a la fitxa d'exercicis)
\newcommand{\recorda}[1]{\par\smallskip\begin{tcolorbox}[colback=accent!7,colframe=accent!55,arc=2mm,boxrule=0.5pt,left=2.5mm,right=2.5mm,top=1mm,bottom=1mm]\small{\color{accent}\faLightbulb[regular]}\ \textbf{Recorda:}\ #1\end{tcolorbox}}
\newcommand{\res}[1]{#1}

\def\fmneed{14}
% =====================================================================
% Capçalera FentMates de les fitxes: un color per a cada curs.
% Els colors són els de la pàgina inicial (LEVELS a assets/js/materials.js);
% el logo ve de _eines/logo/pdf/ (python3 _eines/logo/genera.py).
%
% Ús, al final del preàmbul de la unitat:
%   % =====================================================================
% Capçalera FentMates de les fitxes: un color per a cada curs.
% Els colors són els de la pàgina inicial (LEVELS a assets/js/materials.js);
% el logo ve de _eines/logo/pdf/ (python3 _eines/logo/genera.py).
%
% Ús, al final del preàmbul de la unitat:
%   % =====================================================================
% Capçalera FentMates de les fitxes: un color per a cada curs.
% Els colors són els de la pàgina inicial (LEVELS a assets/js/materials.js);
% el logo ve de _eines/logo/pdf/ (python3 _eines/logo/genera.py).
%
% Ús, al final del preàmbul de la unitat:
%   \input{../comu/fentmates.tex}
%   \fmcurs{1eso}
% i, al començament del document:
%   \capcalera{Matemàtiques}{1r ESO}{Fitxa d'exercicis · Unitat 2}{Divisibilitat}
% Opcional abans de \input: \def\fmneed{14} (espai mínim abans de cada secció).
% =====================================================================
\makeatletter
\def\fm@logodir{../../logo/pdf}
\providecommand{\fmneed}{10}
\def\fm@cursnom{}

% clau del curs, color i color del text damunt del color
\newcommand{\fm@curs}[3]{\@namedef{fm@color@#1}{#2}\@namedef{fm@text@#1}{#3}}
\fm@curs{1eso}{E6B84A}{navy}
\fm@curs{2eso}{3B82F6}{white}
\fm@curs{3eso}{FF7A50}{white}
\fm@curs{4eso-a}{A78BFA}{white}
\fm@curs{4eso-b}{6D28D9}{white}
\fm@curs{1batx-mi}{9B1C1C}{white}
\fm@curs{1batx-ccss}{EF6B6B}{white}
\fm@curs{2batx-mii}{1E6B47}{white}
\fm@curs{2batx-ccss}{2BB3A3}{white}

\newcommand{\fmcurs}[1]{%
  \@ifundefined{fm@color@#1}{\PackageError{fentmates}{Curs desconegut: #1}{}}{}%
  \edef\fm@tmp{\noexpand\definecolor{curs}{HTML}{\@nameuse{fm@color@#1}}%
               \noexpand\colorlet{curstext}{\@nameuse{fm@text@#1}}}\fm@tmp
  \colorlet{cursfosc}{curs!72!black}%
  \def\fm@logo{\fm@logodir/vertical-#1.pdf}%
  \def\fm@marca{\fm@logodir/marca-#1.pdf}}

% ---------- capçalera de la primera pàgina ----------
% #1 matèria (petit), #2 curs, #3 tipus de document i unitat, #4 títol
\newcommand{\capcalera}[4]{%
  \gdef\fm@cursnom{#2}%
  \par\noindent
  \begin{tcolorbox}[colback=curs,colframe=curs,boxrule=0pt,arc=2.2mm,width=3.5cm,height=2.15cm,
      valign=center,left=3mm,right=2mm,top=1mm,bottom=1mm,nobeforeafter,box align=center]
    \color{curstext}{\footnotesize #1}\par\vspace{1pt}{\fontsize{19}{21}\selectfont\bfseries #2\par}%
  \end{tcolorbox}\hspace{5mm}%
  \begin{minipage}[c]{\dimexpr\linewidth-3.5cm-5mm-2.9cm\relax}
    {\color{cursfosc}\small\bfseries\MakeUppercase{#3}}\par\vspace{2pt}
    {\color{navy}\fontsize{20}{23}\selectfont\bfseries #4\par}
  \end{minipage}\hfill
  \begin{minipage}[c]{2.6cm}\raggedleft\includegraphics[width=2.6cm]{\fm@logo}\end{minipage}\par
  \vspace{2pt}}

% ---------- barres de secció del color del curs ----------
\renewcommand{\seccio}[2]{\par\needspace{\fmneed\baselineskip}\bigskip
  \begin{tcolorbox}[colback=curs,colframe=curs,coltext=curstext,arc=2mm,boxrule=0pt,left=3mm,right=3mm,top=1.2mm,bottom=1.2mm]
  {\large\bfseries #1\quad #2}
  \end{tcolorbox}\vspace{-2pt}}

% ---------- capçalera i peu de les pàgines ----------
% primera pàgina: logos institucionals al peu
\fancypagestyle{primera}{\fancyhf{}\renewcommand{\headrulewidth}{0pt}
  \fancyfoot[L]{\raisebox{-0.32\height}{\includegraphics[height=0.78cm]{logo-gva.png}}\hspace{5mm}%
    \raisebox{-0.32\height}{\includegraphics[height=0.78cm]{logo-vermellar.png}}}
  \fancyfoot[R]{\scriptsize\color{muted}\href{https://creativecommons.org/licenses/by-nc-sa/4.0/deed.ca}{\ccbyncsa}\ \ CC BY-NC-SA 4.0 · fentmates.org\qquad\small\thepage}}
% la resta: el símbol a l'esquerra i el curs a la dreta, amb el títol al centre com abans
\fancyhead[L]{\raisebox{-0.3\height}{\includegraphics[height=0.62cm]{\fm@marca}}}
\fancyhead[R]{\small\bfseries\color{cursfosc}\fm@cursnom}
\renewcommand{\headrule}{\color{curs!45}\hrule height \headrulewidth}
\makeatother

%   \fmcurs{1eso}
% i, al començament del document:
%   \capcalera{Matemàtiques}{1r ESO}{Fitxa d'exercicis · Unitat 2}{Divisibilitat}
% Opcional abans de \input: \def\fmneed{14} (espai mínim abans de cada secció).
% =====================================================================
\makeatletter
\def\fm@logodir{../../logo/pdf}
\providecommand{\fmneed}{10}
\def\fm@cursnom{}

% clau del curs, color i color del text damunt del color
\newcommand{\fm@curs}[3]{\@namedef{fm@color@#1}{#2}\@namedef{fm@text@#1}{#3}}
\fm@curs{1eso}{E6B84A}{navy}
\fm@curs{2eso}{3B82F6}{white}
\fm@curs{3eso}{FF7A50}{white}
\fm@curs{4eso-a}{A78BFA}{white}
\fm@curs{4eso-b}{6D28D9}{white}
\fm@curs{1batx-mi}{9B1C1C}{white}
\fm@curs{1batx-ccss}{EF6B6B}{white}
\fm@curs{2batx-mii}{1E6B47}{white}
\fm@curs{2batx-ccss}{2BB3A3}{white}

\newcommand{\fmcurs}[1]{%
  \@ifundefined{fm@color@#1}{\PackageError{fentmates}{Curs desconegut: #1}{}}{}%
  \edef\fm@tmp{\noexpand\definecolor{curs}{HTML}{\@nameuse{fm@color@#1}}%
               \noexpand\colorlet{curstext}{\@nameuse{fm@text@#1}}}\fm@tmp
  \colorlet{cursfosc}{curs!72!black}%
  \def\fm@logo{\fm@logodir/vertical-#1.pdf}%
  \def\fm@marca{\fm@logodir/marca-#1.pdf}}

% ---------- capçalera de la primera pàgina ----------
% #1 matèria (petit), #2 curs, #3 tipus de document i unitat, #4 títol
\newcommand{\capcalera}[4]{%
  \gdef\fm@cursnom{#2}%
  \par\noindent
  \begin{tcolorbox}[colback=curs,colframe=curs,boxrule=0pt,arc=2.2mm,width=3.5cm,height=2.15cm,
      valign=center,left=3mm,right=2mm,top=1mm,bottom=1mm,nobeforeafter,box align=center]
    \color{curstext}{\footnotesize #1}\par\vspace{1pt}{\fontsize{19}{21}\selectfont\bfseries #2\par}%
  \end{tcolorbox}\hspace{5mm}%
  \begin{minipage}[c]{\dimexpr\linewidth-3.5cm-5mm-2.9cm\relax}
    {\color{cursfosc}\small\bfseries\MakeUppercase{#3}}\par\vspace{2pt}
    {\color{navy}\fontsize{20}{23}\selectfont\bfseries #4\par}
  \end{minipage}\hfill
  \begin{minipage}[c]{2.6cm}\raggedleft\includegraphics[width=2.6cm]{\fm@logo}\end{minipage}\par
  \vspace{2pt}}

% ---------- barres de secció del color del curs ----------
\renewcommand{\seccio}[2]{\par\needspace{\fmneed\baselineskip}\bigskip
  \begin{tcolorbox}[colback=curs,colframe=curs,coltext=curstext,arc=2mm,boxrule=0pt,left=3mm,right=3mm,top=1.2mm,bottom=1.2mm]
  {\large\bfseries #1\quad #2}
  \end{tcolorbox}\vspace{-2pt}}

% ---------- capçalera i peu de les pàgines ----------
% primera pàgina: logos institucionals al peu
\fancypagestyle{primera}{\fancyhf{}\renewcommand{\headrulewidth}{0pt}
  \fancyfoot[L]{\raisebox{-0.32\height}{\includegraphics[height=0.78cm]{logo-gva.png}}\hspace{5mm}%
    \raisebox{-0.32\height}{\includegraphics[height=0.78cm]{logo-vermellar.png}}}
  \fancyfoot[R]{\scriptsize\color{muted}\href{https://creativecommons.org/licenses/by-nc-sa/4.0/deed.ca}{\ccbyncsa}\ \ CC BY-NC-SA 4.0 · fentmates.org\qquad\small\thepage}}
% la resta: el símbol a l'esquerra i el curs a la dreta, amb el títol al centre com abans
\fancyhead[L]{\raisebox{-0.3\height}{\includegraphics[height=0.62cm]{\fm@marca}}}
\fancyhead[R]{\small\bfseries\color{cursfosc}\fm@cursnom}
\renewcommand{\headrule}{\color{curs!45}\hrule height \headrulewidth}
\makeatother

%   \fmcurs{1eso}
% i, al començament del document:
%   \capcalera{Matemàtiques}{1r ESO}{Fitxa d'exercicis · Unitat 2}{Divisibilitat}
% Opcional abans de \input: \def\fmneed{14} (espai mínim abans de cada secció).
% =====================================================================
\makeatletter
\def\fm@logodir{../../logo/pdf}
\providecommand{\fmneed}{10}
\def\fm@cursnom{}

% clau del curs, color i color del text damunt del color
\newcommand{\fm@curs}[3]{\@namedef{fm@color@#1}{#2}\@namedef{fm@text@#1}{#3}}
\fm@curs{1eso}{E6B84A}{navy}
\fm@curs{2eso}{3B82F6}{white}
\fm@curs{3eso}{FF7A50}{white}
\fm@curs{4eso-a}{A78BFA}{white}
\fm@curs{4eso-b}{6D28D9}{white}
\fm@curs{1batx-mi}{9B1C1C}{white}
\fm@curs{1batx-ccss}{EF6B6B}{white}
\fm@curs{2batx-mii}{1E6B47}{white}
\fm@curs{2batx-ccss}{2BB3A3}{white}

\newcommand{\fmcurs}[1]{%
  \@ifundefined{fm@color@#1}{\PackageError{fentmates}{Curs desconegut: #1}{}}{}%
  \edef\fm@tmp{\noexpand\definecolor{curs}{HTML}{\@nameuse{fm@color@#1}}%
               \noexpand\colorlet{curstext}{\@nameuse{fm@text@#1}}}\fm@tmp
  \colorlet{cursfosc}{curs!72!black}%
  \def\fm@logo{\fm@logodir/vertical-#1.pdf}%
  \def\fm@marca{\fm@logodir/marca-#1.pdf}}

% ---------- capçalera de la primera pàgina ----------
% #1 matèria (petit), #2 curs, #3 tipus de document i unitat, #4 títol
\newcommand{\capcalera}[4]{%
  \gdef\fm@cursnom{#2}%
  \par\noindent
  \begin{tcolorbox}[colback=curs,colframe=curs,boxrule=0pt,arc=2.2mm,width=3.5cm,height=2.15cm,
      valign=center,left=3mm,right=2mm,top=1mm,bottom=1mm,nobeforeafter,box align=center]
    \color{curstext}{\footnotesize #1}\par\vspace{1pt}{\fontsize{19}{21}\selectfont\bfseries #2\par}%
  \end{tcolorbox}\hspace{5mm}%
  \begin{minipage}[c]{\dimexpr\linewidth-3.5cm-5mm-2.9cm\relax}
    {\color{cursfosc}\small\bfseries\MakeUppercase{#3}}\par\vspace{2pt}
    {\color{navy}\fontsize{20}{23}\selectfont\bfseries #4\par}
  \end{minipage}\hfill
  \begin{minipage}[c]{2.6cm}\raggedleft\includegraphics[width=2.6cm]{\fm@logo}\end{minipage}\par
  \vspace{2pt}}

% ---------- barres de secció del color del curs ----------
\renewcommand{\seccio}[2]{\par\needspace{\fmneed\baselineskip}\bigskip
  \begin{tcolorbox}[colback=curs,colframe=curs,coltext=curstext,arc=2mm,boxrule=0pt,left=3mm,right=3mm,top=1.2mm,bottom=1.2mm]
  {\large\bfseries #1\quad #2}
  \end{tcolorbox}\vspace{-2pt}}

% ---------- capçalera i peu de les pàgines ----------
% primera pàgina: logos institucionals al peu
\fancypagestyle{primera}{\fancyhf{}\renewcommand{\headrulewidth}{0pt}
  \fancyfoot[L]{\raisebox{-0.32\height}{\includegraphics[height=0.78cm]{logo-gva.png}}\hspace{5mm}%
    \raisebox{-0.32\height}{\includegraphics[height=0.78cm]{logo-vermellar.png}}}
  \fancyfoot[R]{\scriptsize\color{muted}\href{https://creativecommons.org/licenses/by-nc-sa/4.0/deed.ca}{\ccbyncsa}\ \ CC BY-NC-SA 4.0 · fentmates.org\qquad\small\thepage}}
% la resta: el símbol a l'esquerra i el curs a la dreta, amb el títol al centre com abans
\fancyhead[L]{\raisebox{-0.3\height}{\includegraphics[height=0.62cm]{\fm@marca}}}
\fancyhead[R]{\small\bfseries\color{cursfosc}\fm@cursnom}
\renewcommand{\headrule}{\color{curs!45}\hrule height \headrulewidth}
\makeatother

\fmcurs{3eso}

\begin{document}
\thispagestyle{primera}

\capcalera{Matemàtiques}{3r ESO}{Fitxa d'exercicis · Unitat 4}{Àlgebra. Polinomis}
\vspace{10pt}
Nom i cognoms: \rule{8.4cm}{0.4pt}\hfill Grup: \rule{1.6cm}{0.4pt}\hfill Data: \rule{2.2cm}{0.4pt}
\vspace{8pt}

\begin{tcolorbox}[colback=soft,colframe=soft,arc=2mm,left=3mm,right=3mm,top=1.5mm,bottom=1.5mm]
\small Exercicis graduats per nivells de dificultat:\quad \estrelles{1}\ bàsic \qquad \estrelles{2}\ mitjà \qquad \estrelles{3}\ avançat.\\[2pt] Resol-los al quadern mostrant tots els passos i dona els polinomis ordenats de major a menor grau. Fes els càlculs sense calculadora.
\end{tcolorbox}

% =====================================================
\seccio{1}{Llenguatge algebraic}

\ex{1} Escriu en llenguatge algebraic:\begin{tasks}(2)
  \task El doble d'un nombre més 5.
  \task El triple d'un nombre menys el seu quadrat.
  \task La meitat del següent d'un nombre.
  \task El quadrat de la suma de dos nombres.
  \task La suma de dos nombres consecutius.
  \task La tercera part d'un nombre més el seu doble.
  \task El $15\,\%$ d'un nombre.
  \task La diferència dels quadrats de dos nombres.
\end{tasks}\recorda{Fixa't en l'ordre: «el quadrat de la suma» és $(x+y)^2$, però «la suma dels quadrats» és $x^2+y^2$.}

\ex{1} Escriu amb paraules el que expressa cada expressió algebraica:\begin{tasks}(4)
  \task $2x-1$
  \task $x^2+3$
  \task $\dfrac{x+y}{2}$
  \task $3(x-4)$
\end{tasks}

\ex{1} Calcula el valor numèric de cada expressió per als valors indicats:\begin{tasks}(2)
  \task $3x-5$, \ per a $x=4$
  \task $x^2-2x$, \ per a $x=-3$
  \task $-x^2+5$, \ per a $x=-2$
  \task $2x^2y-xy$, \ per a $x=1,\ y=-2$
\end{tasks}

\ex{2} Escriu una expressió algebraica per al perímetre i l'àrea de cada figura:\begin{tasks}(2)
  \task Un rectangle de base $3x$ i altura $x+2$.
  \task Un quadrat de costat $x-1$.
\end{tasks}

\ex{2} En Pere té $x$ anys; la seua germana en té el doble menys 3 i son pare, 25 més que en Pere.\begin{tasks}(1)
  \task Escriu l'edat de cadascú i la suma de les tres edats.
  \task Escriu la suma de les tres edats d'ací a 4 anys.
  \task Quant sumen les tres edats si en Pere té 12 anys?
\end{tasks}

\ex{2} Un aparcament cobra 1,80 € per entrar i 0,05 € per cada minut. Escriu l'expressió del preu d'una estada de $x$ minuts i calcula quant costen 2 hores.

\ex{3} Fes servir el llenguatge algebraic per a demostrar que:\begin{tasks}(1)
  \task Si a un nombre li sumem el seu doble, el resultat sempre és múltiple de 3.
  \task La suma de dos nombres parells sempre és parell.
  \task La suma de dos nombres senars consecutius sempre és múltiple de 4.
\end{tasks}


% =====================================================
\seccio{2}{Monomis}

\ex{1} Copia i completa la taula:\par\smallskip{\centering\renewcommand{\arraystretch}{1.35}\begin{tabular}{|C{2.6cm}|C{2.6cm}|C{2.6cm}|C{1.6cm}|}\hline\rowcolor{soft}\textbf{Monomi} & \textbf{Coeficient} & \textbf{Part literal} & \textbf{Grau}\\\hline $-8x^3$ & & & \\\hline $a^2b$ & & & \\\hline $\tfrac{2}{3}xy^4$ & & & \\\hline $-y$ & & & \\\hline $12$ & & & \\\hline $5x^2y^3z$ & & & \\\hline \end{tabular}\par}\recorda{El \textbf{grau} d'un monomi és la suma dels exponents de la part literal. Si no s'escriu el coeficient, val 1.}

\ex{1} Quines d'aquestes expressions són monomis? Justifica la resposta.\begin{tasks}(4)
  \task $-3x^4$
  \task $\dfrac{5}{x}$
  \task $\dfrac{x^2}{4}$
  \task $2x+y$
  \task $\sqrt{x}$
  \task $7$
  \task $x^{-2}y$
  \task $-\tfrac12ab^2$
\end{tasks}

\ex{1} Agrupa els monomis semblants:\quad $4x^2,\ \ -xy,\ \ 3x,\ \ -\tfrac12x^2,\ \ 5yx,\ \ x,\ \ 2x^2y,\ \ -x^2y,\ \ 9$.

\ex{2} Indica el grau de cada monomi respecte de cada variable i el grau total. Escriu, per a cadascun, un monomi semblant.\begin{tasks}(3)
  \task $-6x^2y^5$
  \task $ab^3c$
  \task $\tfrac{3}{5}x^4y$
\end{tasks}

\ex{2} Calcula el valor numèric:\begin{tasks}(3)
  \task $-4x^3$ \ per a $x=-2$
  \task $3xy^2$ \ per a $x=2,\ y=-1$
  \task $\tfrac12x^2y$ \ per a $x=4,\ y=3$
\end{tasks}

\ex{3} Resol:\begin{tasks}(1)
  \task Troba $k$ perquè el monomi $-2x^{k}y^3$ siga de grau 8.
  \task Troba $m$ i $n$ perquè $5x^{m+1}y^2$ i $-3x^4y^{n-1}$ siguen semblants.
  \task Escriu un monomi de grau 6 amb tres variables, coeficient $-\tfrac23$ i grau 3 respecte de $x$.
\end{tasks}


% =====================================================
\seccio{3}{Operacions amb monomis}

\ex{1} Calcula:\begin{tasks}(3)
  \task $3x+5x$
  \task $7x^{2}-2x^{2}$
  \task $-4ab+ab$
  \task $2x^{3}-x^{3}+5x^{3}$
  \task $xy-3xy-2xy$
  \task $6a^{2}b-6a^{2}b$
\end{tasks}\recorda{Només es poden sumar i restar monomis \textbf{semblants}: sumem els coeficients i \textbf{no canviem} la part literal.}

\ex{1} Redueix els termes semblants:\begin{tasks}(2)
  \task $3x^{2}-5x+x^{2}+2x$
  \task $4xy-2x+3xy+7x$
  \task $2a^{2}-a+5-a^{2}+3a-8$
  \task $x^{3}-2x+3x^{3}+x^{2}-4x$
\end{tasks}

\ex{1} Multiplica:\begin{tasks}(3)
  \task $2x^{3}\cdot 5x^{2}$
  \task $(-3a^{2})\cdot (4a^{3})$
  \task $(6xy^{2})\cdot (-2x^{3}y)$
  \task $(\tfrac{1}{2}x^{4})\cdot (8x)$
  \task $(-x^{2}y)\cdot (-xy^{3})\cdot 2x$
  \task $(-5ab)\cdot (3a^{2}b^{2})$
\end{tasks}\recorda{Per a multiplicar, multipliquem els coeficients i \textbf{sumem els exponents} de cada lletra: $x^m\cdot x^n=x^{m+n}$.}

\ex{2} Divideix. Indica si el resultat és un monomi.\begin{tasks}(3)
  \task $12x^5:(3x^2)$
  \task $(-20a^4b^3):(5a^2b)$
  \task $9x^3y^2:(-3x^3y)$
  \task $(-14x^6):(-7x^6)$
  \task $6x^2y:(4x)$
  \task $4x^2:(8x^3)$
\end{tasks}

\ex{2} Calcula les potències:\begin{tasks}(4)
  \task $(2x^{3})^{2}$
  \task $(-3xy^{2})^{3}$
  \task $(-x^{2}y)^{4}$
  \task $(\tfrac{1}{2} x)^{2}$
\end{tasks}

\ex{2} Opera i simplifica (recorda la jerarquia):\begin{tasks}(2)
  \task $3x\cdot2x^2-4x^3$
  \task $10x^5:(2x^2)+x\cdot x^2$
  \task $(2x)^3-5x\cdot x^2$
  \task $6a^2b:(3a)-ab\cdot4$
\end{tasks}

\ex{3} Troba el monomi que falta en cada cas:\begin{tasks}(1)
  \task $(\quad?\quad)\cdot3x^2=12x^5$
  \task $20x^4y^2:(\quad?\quad)=-5x^2y$
  \task $4x^3+(\quad?\quad)=-x^3$
  \task $(\quad?\quad)^2=9x^6y^2$ \ (dues solucions)
\end{tasks}


% =====================================================
\seccio{4}{Polinomis}

\ex{1} Per a cada polinomi: ordena'l, i indica'n el grau, el nombre de termes, el coeficient principal i el terme independent.\begin{tasks}(2)
  \task $P(x)=5x-3x^{4}+2-x^{2}$
  \task $P(x)=7-x^{3}$
  \task $P(x)=4x^{2}-3x+x^{5}$
  \task $P(x)=-2+x$
\end{tasks}\recorda{El \textbf{grau} d'un polinomi és el grau del terme de grau més alt. El \textbf{terme independent} és el que no té lletra.}

\ex{1} Quins polinomis de l'exercici anterior són complets? Escriu els incomplets en forma completa (amb coeficients 0).

\ex{1} Siga $P(x)=2x^3-x^2+3x-4$. Calcula:\begin{tasks}(4)
  \task $P(0)$
  \task $P(1)$
  \task $P(-2)$
  \task $P\!\left(\tfrac12\right)$
\end{tasks}

\ex{2} Indica el grau de cada polinomi:\begin{tasks}(3)
  \task $3x^2y^3-xy+5$
  \task $a^4-2a^2b^3+b$
  \task $x^2y^2-x^3+7xy$
\end{tasks}

\ex{2} Comprova quins dels nombres $1,\ -1,\ 2,\ -2,\ 3$ i $-3$ són arrels del polinomi $P(x)=x^{3}-7x+6$.

\ex{3} Resol:\begin{tasks}(1)
  \task Troba $m$ perquè $P(x)=2x^3-mx+4$ complisca $P(-2)=0$.
  \task Troba $k$ perquè el polinomi $Q(x)=x^2+kx-6$ tinga per arrel $x=3$.
  \task Escriu un polinomi de grau 4, incomplet, amb coeficient principal $-1$ i tal que $P(0)=5$.
\end{tasks}


% =====================================================
\seccio{5}{Suma i resta de polinomis}

\ex{1} Siguen els polinomis $P(x)=3x^{3}-2x^{2}+x-5$, \ $Q(x)=-x^{3}+4x^{2}-6$ \ i \ $R(x)=2x^{2}-3x+1$. Calcula:\begin{tasks}(4)
  \task $P+Q$
  \task $P-Q$
  \task $Q+R$
  \task $R-P$
\end{tasks}\recorda{Per a restar, sumem l'\textbf{oposat}: canviem el signe de \textbf{tots} els termes del polinomi que restem.}

\ex{1} Fes aquestes sumes i restes col·locant els termes semblants en columna:\begin{tasks}(1)
  \task $(4x^4-3x^2+2x)+(x^3+5x^2-2x+7)$
  \task $(x^3-6x+1)-(2x^3+x^2-6x-4)$
\end{tasks}

\ex{2} Amb els polinomis de l'exercici 27, calcula:\begin{tasks}(3)
  \task $P-Q+R$
  \task $P-(Q-R)$
  \task $P-(Q+R)$
\end{tasks}

\ex{2} Amb els polinomis de l'exercici 27, calcula $2P+3Q$ i $P-2R$.

\ex{2} Opera i redueix:\begin{tasks}(1)
  \task $(5x^2y-3xy+y^2)-(2x^2y+xy-4y^2)$
  \task $(a^2-2ab+b^2)+(a^2+2ab+b^2)-(2a^2-b^2)$
\end{tasks}

\ex{3} Troba el polinomi $X(x)$ en cada cas:\begin{tasks}(1)
  \task $P(x)+X(x)=Q(x)$, \ amb $P$ i $Q$ els de l\'exercici 27.
  \task $(2x^3+x-1)-X(x)=x^3-4x^2+3$
\end{tasks}


% =====================================================
\seccio{6}{Producte de polinomis}

\ex{1} Multiplica:\begin{tasks}(2)
  \task $3x(2x^{2}-5x+4)$
  \task $-2x^{2}(x^{3}-3x+1)$
  \task $4ab(a-2b)$
  \task $-x(x^{2}-x-1)$
\end{tasks}\recorda{Multipliquem el monomi per \textbf{cada terme} del polinomi (propietat distributiva).}

\ex{1} Multiplica:\begin{tasks}(2)
  \task $(x+3)(x+5)$
  \task $(x-4)(x+2)$
  \task $(2x-1)(x-3)$
  \task $(3x+2)(2x-5)$
\end{tasks}

\ex{2} Multiplica:\begin{tasks}(1)
  \task $(x^{2}-3x+2)(x+4)$
  \task $(2x^{2}+x-1)(x-3)$
  \task $(x^{3}-2x+1)(x^{2}+x)$
  \task $(x^{2}+x+1)(x^{2}-x+1)$
\end{tasks}

\ex{2} Siguen $P(x)=x^2-2x+3$ i $Q(x)=2x-1$. Calcula $P\cdot Q$ i $Q\cdot Q$. Quin grau té cada producte?

\ex{2} Opera i redueix:\begin{tasks}(1)
  \task $2x(x-3)-(x+1)(x-2)$
  \task $(x+4)(x-1)-x(x+3)$
  \task $3(x^2-1)-(2x+1)(x-2)$
\end{tasks}

\ex{3} Calcula $(x-1)(x+2)(x+3)$ i $(x+1)(x-1)(x^2+1)$.

\ex{3} D'una planxa rectangular de $(2x+3)$ cm per $(x+4)$ cm en retallem un quadrat de costat $x$ cm. Escriu com a polinomi l'àrea que queda i calcula-la per a $x=5$.


% =====================================================
\seccio{7}{Identitats notables}

\ex{1} Desenvolupa:\begin{tasks}(3)
  \task $(x+3)^{2}$
  \task $(x-5)^{2}$
  \task $(x+4)(x-4)$
  \task $(a+1)^{2}$
  \task $(x-1)^{2}$
  \task $(y+7)(y-7)$
\end{tasks}\recorda{$(a+b)^2=a^2+2ab+b^2$, \quad $(a-b)^2=a^2-2ab+b^2$, \quad $(a+b)(a-b)=a^2-b^2$.}

\ex{2} Desenvolupa:\begin{tasks}(3)
  \task $(2x+5)^{2}$
  \task $(3x-1)^{2}$
  \task $(2x+3y)(2x-3y)$
  \task $(x^{2}-2)^{2}$
  \task $(\tfrac{1}{2}x+2)^{2}$
  \task $(-x-3)^{2}$
\end{tasks}

\ex{2} Expressa com a quadrat d'un binomi o com a suma per diferència:\begin{tasks}(3)
  \task $x^2+6x+9$
  \task $x^2-10x+25$
  \task $x^2-81$
  \task $4x^2+4x+1$
  \task $9x^2-16$
  \task $25x^2-30x+9$
\end{tasks}

\ex{2} Completa:\begin{tasks}(1)
  \task $(x+\ \square\ )^2=x^2+\ \square\ +49$
  \task $(\ \square\ -3)^2=4x^2-\ \square\ +9$
  \task $(x+\ \square\ )(x-\ \square\ )=x^2-64$
  \task $(3x+\ \square\ )^2=\ \square\ +12x+\ \square$
\end{tasks}

\ex{2} Opera i simplifica:\begin{tasks}(1)
  \task $(x+2)^2-(x-2)^2$
  \task $(x+3)(x-3)-(x-3)^2$
  \task $2(x+1)^2-(x+1)(x-1)$
\end{tasks}

\ex{3} Calcula mentalment amb les identitats notables:\begin{tasks}(4)
  \task $101^2$
  \task $99^2$
  \task $51\cdot49$
  \task $1002^2$
\end{tasks}

\ex{3} Demostra que la diferència entre els quadrats de dos nombres consecutius sempre és un nombre senar.


% =====================================================
\seccio{8}{Extracció de factor comú}

\ex{1} Trau factor comú:\begin{tasks}(3)
  \task $5x^{2}-15x$
  \task $8x^{3}+12x^{2}$
  \task $6x^{4}-9x^{3}+3x^{2}$
  \task $7x-14$
  \task $x^{5}-x^{3}+x^{2}$
  \task $10x^{3}-4x$
\end{tasks}\recorda{Factor comú: el MCD dels coeficients i les lletres que estan en \textbf{tots} els termes amb el \textbf{menor} exponent. Comprova-ho multiplicant!}

\ex{1} Trau factor comú:\begin{tasks}(2)
  \task $6x^{2}y-9xy^{2}$
  \task $4a^{3}b^{2}+2a^{2}b$
  \task $12x^{2}y^{3}-8xy^{2}+4xy$
  \task $15a^{2}-10ab+5a$
\end{tasks}

\ex{2} Trau factor comú un binomi:\begin{tasks}(3)
  \task $x(x-2)+5(x-2)$
  \task $3(x+1)-x(x+1)$
  \task $2x(x+4)-(x+4)$
\end{tasks}

\ex{2} Factoritza al màxim: primer trau factor comú i després aplica una identitat notable.\begin{tasks}(3)
  \task $x^{3}-4x$
  \task $2x^{2}-8x+8$
  \task $3x^{3}+6x^{2}+3x$
  \task $x^{4}-16x^{2}$
  \task $5x^{2}-45$
  \task $x^{3}-x$
\end{tasks}

\ex{3} Calcula de la manera més ràpida traient factor comú:\begin{tasks}(3)
  \task $13\cdot47+13\cdot53$
  \task $98^2+2\cdot98$
  \task $37\cdot25+37\cdot75$
\end{tasks}

\ex{3} A un rectangle de costats $x$ i $x+3$ li llevem un rectangle de costats $3$ i $x+3$. Escriu l'àrea que queda, trau factor comú i comprova que és una diferència de quadrats.


% =====================================================
\seccio{9}{Divisió de polinomis}

\ex{1} Fes aquestes divisions i indica el quocient i el residu:\begin{tasks}(3)
  \task $(x^{2}+7x+10):(x+2)$
  \task $(x^{2}-3x+5):(x-1)$
  \task $(2x^{2}+x-3):(x-1)$
\end{tasks}\recorda{Seguim dividint fins que el grau del residu siga \textbf{menor} que el del divisor. Comprovació: $D=d\cdot C+R$.}

\ex{1} Divideix. Compte amb els termes que falten en el dividend:\begin{tasks}(2)
  \task $(x^{3}-2x^{2}+3x-1):(x-2)$
  \task $(3x^{3}+x-4):(x+1)$
\end{tasks}

\ex{2} Divideix:\begin{tasks}(1)
  \task $(x^{4}-2x^{3}+x-3):(x^{2}+1)$
  \task $(2x^{4}+x^{3}-3x^{2}+x):(x^{2}-x+2)$
\end{tasks}

\ex{2} Divideix $(6x^3-x^2+4x-5):(2x-1)$ i comprova el resultat amb la prova de la divisió.

\ex{2} En una divisió de polinomis, el divisor és $x^2-3$, el quocient $2x+1$ i el residu $x-4$. Quin és el dividend?

\ex{3} En dividir $D(x)=x^3+x^2-3x+5$ entre un polinomi $d(x)$ obtenim de quocient $x+2$ i de residu $-2x+3$. Troba $d(x)$.


% =====================================================
\seccio{10}{Regla de Ruffini}

\ex{1} Divideix aplicant la regla de Ruffini:\begin{tasks}(2)
  \task $(x^{3}-3x^{2}+2x-4):(x-2)$
  \task $(x^{2}+6x+5):(x+1)$
  \task $(2x^{3}-5x^{2}+x+3):(x-3)$
\end{tasks}\recorda{Escriu els coeficients del dividend \textbf{ordenat i complet}. Si el divisor és $x-a$, a l'esquerra posem $a$ (per a $x+2$ posem $-2$).}

\ex{1} Divideix per Ruffini. Compte amb els termes que falten:\begin{tasks}(2)
  \task $(x^{4}-3x^{2}+2x-1):(x+2)$
  \task $(3x^{3}-4x+5):(x-1)$
\end{tasks}

\ex{2} Divideix per Ruffini:\begin{tasks}(2)
  \task $(x^{4}-16):(x-2)$
  \task $(x^{5}+1):(x+1)$
\end{tasks}

\ex{2} Calcula, sense fer la divisió, el residu de:\begin{tasks}(3)
  \task $(x^3-5x+7):(x+3)$
  \task $(2x^4-x^2+1):(x-1)$
  \task $(x^{5}-2x^3+x):(x+2)$
\end{tasks}

\ex{2} Comprova, amb el teorema del residu, quins dels nombres $1,\ -1,\ 2,\ -2,\ 3,\ -3$ són arrels de $P(x)=x^{3}-3x^{2}-4x+12$.

\ex{3} Resol:\begin{tasks}(1)
  \task Troba $k$ perquè la divisió $(2x^3+kx^2-x+3):(x+1)$ siga exacta.
  \task Troba $m$ perquè el residu de $(x^3-mx+4):(x-2)$ siga 4.
\end{tasks}

\ex{3} Comprova que $x=1$ és una arrel de $P(x)=x^3-2x^2-x+2$. Divideix per Ruffini entre $x-1$ i, aplicant Ruffini una altra vegada al quocient, escriu $P(x)$ com a producte de tres factors de grau 1.

\vfill
\begin{center}\color{muted}\small\faCheckCircle[regular]\ \ Ànim! Si et costa algun exercici, repassa els apunts interactius de la unitat.\end{center}

\end{document}
